Notre recherche répond à une difficulté d'accès aux connaissances scientifiques sur les plantes médicinales du Burkina Faso. Ces connaissances sont dispersées entre des articles, des thèses, des mémoires et des rapports dont les terminologies et les noms de plantes varient. Il est donc difficile de retrouver une information sur une plante, une pathologie ou une recette et de remonter à sa source. Pour répondre à cette difficulté, nous avons adapté le modèle conceptuel d'une ontologie existante, enrichi l'ontologie en données issues de la littérature scientifique, puis importé et ajusté cette représentation dans une base orientée graphe exploitée par un système d'accès intelligent et une application web. La démarche a suivi six phases, de la constitution du corpus à l'évaluation du dispositif, avec deux volets menés en parallèle sur le modèle conceptuel et sur le modèle de données.

Nous avons comparé les ontologies du domaine au moyen d'une grille commune et retenu \textit{ontoMEDTRAD}, qui a obtenu le score le plus élevé devant \textit{WATRIMed}, comme base conceptuelle. En phase 3a, nous avons adapté son modèle conceptuel au contexte burkinabè ; en phase 3b, nous avons extrait les connaissances de dix documents scientifiques, puis, en phase 4, enrichi l'ontologie avec ces individus associés à leurs sources. Au cours de la même phase, nous avons importé dans Neo4j l'ontologie réajustée \textit{OntoPlantMedBF}, afin d'exploiter son modèle et ses données dans la base graphe. La comparaison des schémas montre 40 classes ajoutées, 20 propriétés objet ajoutées et 15 propriétés de données ajoutées par rapport au schéma conservé. Ces ajouts portent notamment sur les noms vernaculaires, les usages et maladies cibles, les parties et préparations, les niveaux de preuve, la provenance, la localisation, les composés chimiques, les activités pharmacologiques et les effets toxicologiques rapportés. Nous avons mis en place un système RAG associant Neo4j et Llama 3 8B, exécuté sur Google Colab. Le modèle traduit les questions en langage naturel en requêtes Cypher, validées en lecture seule, et produit des réponses à partir des faits retrouvés. Une application web fournit une interface de recherche et de consultation. Nous avons évalué l'ontologie et le système RAG ; l'application n'a pas été évaluée auprès de ses utilisateurs.

Nous avons examiné les hypothèses à partir des résultats présentés aux chapitres \ref{chap:conception_developpement} et \ref{chap:resultats_discussion}. Faute de campagne de tests auprès des utilisateurs, nos conclusions sur l'application web se limitent à ses fonctions et à ses temps de réponse.

\textbf{Hypothèse 1.} Une revue méthodique des ontologies existantes devait permettre de sélectionner un modèle capable, après réajustement, de représenter les connaissances sur les plantes médicinales du Burkina Faso. Les travaux réalisés étayent cette hypothèse. Nous avons comparé plusieurs modèles, puis construit et peuplé OntoPlantMedBF à partir de l'ontologie retenue, afin de représenter les plantes, leurs usages, leurs parties utilisées, leurs préparations, leurs propriétés et leurs sources. HermiT n'a signalé aucune incohérence dans les vérifications rapportées.

\textbf{Hypothèse 2.} L'intelligence artificielle, associée à une vérification manuelle, devait permettre d'extraire de la littérature scientifique des connaissances fidèles aux documents et traçables jusqu'à leurs sources. Les résultats du tableau \ref{tab:res_ontologie} étayent cette hypothèse : les 289 plantes et les 1~676 usages sont reliés à une source, sur dix des 35 documents du corpus, et le nom de la plante figure dans le document d'origine pour 1~497 usages sur 1~501 contrôlables. La présence de ces liens ne garantit pas l'exactitude des faits, contrôlée seulement sur une partie des documents. Nous n'avons pas non plus comparé cette méthode à une extraction manuelle dans des conditions contrôlées ; nous ne pouvons donc pas conclure qu'elle lui soit supérieure.

\textbf{Hypothèse 3.} Une solution numérique interrogeable en langage naturel devait faciliter l'accès aux connaissances documentées. Nous avons testé le système RAG sur un jeu de 50 questions : avec Llama 3 8B, le F1-score atteint 0,728, 36 réponses sont conformes, et les 5 questions hors corpus reçoivent un refus. Le système restitue donc en partie le contenu du graphe et refuse à tort 6 questions, tandis que pour 2 autres le modèle affirme à tort n'avoir rien trouvé ; Gemini atteint 0,995 sur le même jeu. Ces mesures ne renseignent pas sur la justesse scientifique des faits. Le jeu de questions provient en outre du corpus utilisé pour construire le système, et l'évaluation a été menée par ses concepteurs. L'application a été développée, mais elle n'a pas été évaluée auprès de ses utilisateurs : son apport pour la consultation des connaissances reste à mesurer. Avec le modèle ouvert retenu, l'hypothèse n'est donc que partiellement étayée.

Six limites encadrent ce bilan. L'évaluation a été conduite par nous-mêmes, les concepteurs du système, sans regard extérieur. Le jeu de questions provient du même corpus que celui utilisé pour construire la base et mesure donc surtout les capacités du système sur ces données. Les niveaux de preuve sont repris des sources sans réévaluation indépendante, et les extractions n'ont pas été relues par un expert. Le système reste un prototype déployé localement, dont l'utilisation à grande échelle n'a pas été étudiée. Enfin, il s'agit d'un outil de consultation. Il ne pose pas de diagnostic et n'établit pas de prescription ; les résultats ne constituent pas une validation clinique.

Ces limites indiquent plusieurs prolongements possibles : faire évaluer les extractions par un spécialiste, tester l'application auprès d'utilisateurs, dont des tradipraticiens, élargir le corpus aux thèses et rapports non numérisés, puis mesurer l'effet de ces ajouts sur la couverture. Une relecture experte d'un échantillon et des questions formulées par des utilisateurs compléteraient l'évaluation. Un déploiement pérenne et une collecte discrète des données d'usage permettraient aussi d'étudier le système en conditions réelles.

L'apport du travail tient à l'organisation des connaissances et au lien établi entre les sources, l'ontologie, le graphe et les réponses présentées à l'utilisateur. La base reste de taille modeste, et l'évaluation par les utilisateurs reste à mener. La traçabilité est vérifiée dans le graphe ; des contrôles dans les documents et des évaluations plus larges seront nécessaires pour en établir la portée.
